\subsection{Rule-Shift Adaptation}
\label{sec:rule_shift}

\paragraph{Task Description.}
This experiment injects a mid-stream rule change into the substrate without
any explicit signal or re-training.  Two generative rules partition the
$2N$ = 50-sample stream:

\begin{itemize}[nosep]
  \item \textbf{Rule A} (steps 0--24 inclusive, $N=25$ samples):
        linear modular sequences, $b_k = (s_i + 3k) \bmod 251$.
  \item \textbf{Rule B} (steps 25--49 inclusive, $N=25$ samples):
        XOR-nonlinear sequences, $b_k = s_i \oplus (11k \bmod 256)$.
\end{itemize}

At every step the pipeline's retrieval output is scored against the
expected continuation under \emph{both} rules simultaneously, yielding
per-step alignment signals $K_A$ and $K_B$.  The winner at each step
is $\arg\max(K_A, K_B)$.

\paragraph{Results.}
Figure~\ref{fig:ruleshift_trajectory} shows the full adaptation
trajectory.  The shift boundary is at step 25 (the first index governed by Rule B).  Rule B
did not stabilise within the observation window.  Phase A had $23$ of
$25$ steps won by Rule A
(92\,\%); Phase B had $0$ of
$25$ steps won by Rule B (0\,\%).

\paragraph{Assessment.}
The substrate did not complete the transition within the available window.
At this scale, Rules A and B produce similar value cosine similarities;
longer samples or larger \textit{N} would widen the attractor gap and make
the shift detectable.
